% GENERATED FILE. Edit canonical lessons, then run npm run generate:books.
% canonical-source-hash: fnv1a64:835ad82a56a06d5d
% canonical-lessons: LA-C279-cerasum, LA-C279-prunum, LA-C279-brassica, LA-C279-porrum, LA-C279-puls

\chapter{Cherry, Plum, Cabbage, Leek, Porridge}
\label{ch:la-cherry-plum-cabbage-leek-porridge}
\hlchaptermodality{279}

\begin{chapteropening}
\textbf{By the end of this chapter:} \emph{I can say a cherry, a plum, cabbage, a leek and porridge.}

Complete the last lesson of chapter 279: I can say a cherry, a plum, cabbage, a leek and porridge.
\end{chapteropening}

\section[cerasum, cerasī]{cerasum, cerasī — a cherry}

\label{lesson:LA-C279-cerasum}



\begin{quote}
Before the new one: say the Latin for a word, then the Latin for a name.
\end{quote}

\subsection*{You'll want to know: cerasum, cerasī}
\textbf{cerasum, cerasī} — \textquotedblleft{}a cherry\textquotedblright{}.

\begin{grammarlens}[title={What you've built}]
One more word for this chapter.
\end{grammarlens}

\subsection*{Your turn}
\begin{itemize}
\raggedright
  \item \emph{Say it:} \emph{cerasum, cerasī}
  \item \emph{Say it:} \emph{cerasum, cerasī}, once more
  \item \emph{Say it:} say \emph{cerasum, cerasī}
  \item \emph{Recall:} say the Latin for a word, then the Latin for a name, then say \emph{cerasum, cerasī} again
\end{itemize}

\begin{tcolorbox}[breakable,colback=teal!4,colframe=teal!35!black,title={Before you move on}]
What does cerasum, cerasī mean? (\textquotedblleft{}A cherry\textquotedblright{}.) Say it once more.
\end{tcolorbox}
\section[prūnum, prūnī]{prūnum, prūnī — a plum}

\label{lesson:LA-C279-prunum}



\begin{quote}
Before the new one: say the Latin for a name, then the Latin for a cherry.
\end{quote}

\subsection*{You'll want to know: prūnum, prūnī}
\textbf{prūnum, prūnī} — \textquotedblleft{}a plum\textquotedblright{}.

\begin{grammarlens}[title={What you've built}]
One more word for this chapter.
\end{grammarlens}

\subsection*{Your turn}
\begin{itemize}
\raggedright
  \item \emph{Say it:} \emph{prūnum, prūnī}
  \item \emph{Say it:} \emph{prūnum, prūnī}, once more
  \item \emph{Say it:} say \emph{prūnum, prūnī}
  \item \emph{Recall:} say the Latin for a name, then the Latin for a cherry, then say \emph{prūnum, prūnī} again
\end{itemize}

\begin{tcolorbox}[breakable,colback=teal!4,colframe=teal!35!black,title={Before you move on}]
What does prūnum, prūnī mean? (\textquotedblleft{}A plum\textquotedblright{}.) Say it once more.
\end{tcolorbox}
\section[brassica, brassicae]{brassica, brassicae — cabbage}

\label{lesson:LA-C279-brassica}



\begin{quote}
Before the new one: say the Latin for a cherry, then the Latin for a plum.
\end{quote}

\subsection*{You'll want to know: brassica, brassicae}
\textbf{brassica, brassicae} — \textquotedblleft{}cabbage\textquotedblright{}.

\begin{grammarlens}[title={What you've built}]
One more word for this chapter.
\end{grammarlens}

\subsection*{Your turn}
\begin{itemize}
\raggedright
  \item \emph{Say it:} \emph{brassica, brassicae}
  \item \emph{Say it:} \emph{brassica, brassicae}, once more
  \item \emph{Say it:} say \emph{brassica, brassicae}
  \item \emph{Recall:} say the Latin for a cherry, then the Latin for a plum, then say \emph{brassica, brassicae} again
\end{itemize}

\begin{tcolorbox}[breakable,colback=teal!4,colframe=teal!35!black,title={Before you move on}]
What does brassica, brassicae mean? (\textquotedblleft{}Cabbage\textquotedblright{}.) Say it once more.
\end{tcolorbox}
\section[porrum, porrī]{porrum, porrī — a leek}

\label{lesson:LA-C279-porrum}



\begin{quote}
Before the new one: say the Latin for a plum, then the Latin for cabbage.
\end{quote}

\subsection*{You'll want to know: porrum, porrī}
\textbf{porrum, porrī} — \textquotedblleft{}a leek\textquotedblright{}.

\begin{grammarlens}[title={What you've built}]
One more word for this chapter.
\end{grammarlens}

\subsection*{Your turn}
\begin{itemize}
\raggedright
  \item \emph{Say it:} \emph{porrum, porrī}
  \item \emph{Say it:} \emph{porrum, porrī}, once more
  \item \emph{Say it:} say \emph{porrum, porrī}
  \item \emph{Recall:} say the Latin for a plum, then the Latin for cabbage, then say \emph{porrum, porrī} again
\end{itemize}

\begin{tcolorbox}[breakable,colback=teal!4,colframe=teal!35!black,title={Before you move on}]
What does porrum, porrī mean? (\textquotedblleft{}A leek\textquotedblright{}.) Say it once more.
\end{tcolorbox}
\section[puls, pultis]{puls, pultis — porridge}

\label{lesson:LA-C279-puls}



\begin{quote}
Before the new one: say the Latin for cabbage, then the Latin for a leek.
\end{quote}

\subsection*{You'll want to know: puls, pultis}
\textbf{puls, pultis} — \textquotedblleft{}porridge\textquotedblright{}.

\begin{grammarlens}[title={What you've built}]
One more word for this chapter.
\end{grammarlens}

\subsection*{Your turn}
\begin{itemize}
\raggedright
  \item \emph{Say it:} \emph{puls, pultis}
  \item \emph{Say it:} \emph{puls, pultis}, once more
  \item \emph{Say it:} say \emph{puls, pultis}
  \item \emph{Recall:} say the Latin for cabbage, then the Latin for a leek, then say \emph{puls, pultis} again
\end{itemize}

\begin{tcolorbox}[breakable,colback=teal!4,colframe=teal!35!black,title={Before you move on}]
What does puls, pultis mean? (\textquotedblleft{}Porridge\textquotedblright{}.) Say it once more.
\end{tcolorbox}
